\section{Del 1: DNBs argumenter i lys av Teori X og Teori Y}

\subsection{DNBs beslutning og teoriens premisser}

DNB har gjort innstramminger i regelverket for hjemmekontor; 

\enquote{Idag kan ansatte i DNB jobbe hjemmefra inntil to dager i uken. Fra oktober endres ordningen: Kontoret blir utgangspunktet, mens hjemmekontor må avklares med nærmeste leder} \parencite{biernat2026}.

DNB argumenterer for at innstrammingene vil bedre arbeidsmiljøet, samt skape flere muligheter for læring og samarbeid. Likevel har Innstrammingene skapt sterke reaksjoner hos flere av de ansatte i bedriften \parencite{biernat2026}. DNB svarer på disse reaksjonene ved å understreke at autonomien til de ansatte ikke er fullstendig borte fordi det fortsatt er mulig å avtale hjemmekontor med nærmeste leder \parencite{biernat2026}. Ved hjelp av Douglas McGregor sin Teori X og teori Y \parencite[s.~427]{einarsen2023} vil denne delen av oppgaven forsøke å drøfte legitimiteten til DNB sine argumenter rundt endringene i regelverket. Og forsøke å sette søkelys på de ulike utfallene og konsekvensene slike endringer i regelverk kan medføre.

\subsection{Innstrammingen vurdert gjennom Teori X}

Douglas McGregor har skapt en teori der han presenterer to ulike forskjellige typer menneskesyn. Denne teorien heter teori X og teori Y. McGregors teori X går ut på at; \enquote{mennesker er grunnleggende late, og at de alltid vil unngå å jobbe så sant det er mulig, samt at de kun motiveres av lønn eller fritid} \parencite[s.~427]{einarsen2023}.

Det er mulig å anvende menneskesynet presentert i teori X for å argumentere for at DNB sine innstramminger vil fungere godt. Dersom mennesket motiveres av lønn og fritid, og derfor ikke ønsker å arbeide uten ytre motivasjon, vil det være usikkert om arbeidstakere utnytter sitt fulle potensiale på hjemmekontor.

DNB argumenterer for at mindre hjemmekontor vil skape et bedre arbeidsmiljø \parencite{biernat2026}. Ettersom teori X forteller at lønn og fritid er det eneste som motiverer mennesker til å arbeide, kan det være vanskelig å arbeide i omgivelser som minner om fritid. Derfor vil DNB sine argumenter om at; \enquote{Det styrker arbeidsmiljøet vårt, gir bedre muligheter for læring og samarbeid} \parencite{biernat2026} fungere godt. Det er fordi fritiden til de ansatte foregår utenfor arbeidsplassen og sannsynligvis på liknende arenaer som hjemmekontoret foregår. Disse likhetene mellom fritid og hjemmekontor kan være argumentere for at arbeidsplassen er et bedre sted for gode arbeidsvilkår og godt læringsmiljø. DNB argumenterer også slik; \enquote{-- Kontoret må oppleves som et sted det er meningsfullt å være, fordi man løser oppgaver, samarbeider, lærer og hører til.} \parencite{biernat2026}.

\subsection{Kritikk av Teori X og Argyris' perspektiv på kontroll}

På den andre siden er det mulig å bruke teori X for å forklare hvorfor DNB sine argumenter ikke fungerer godt. DNB argumenterer for at Innstramminger rundt hjemmekontorordningen vil føre til bedre arbeidsmiljø og styrke mulighetene for læring og samarbeid. McGregor sin teori på sin side argumenterer at dersom ledere baserer sin lederstil på et teori X menneskesyn, vil det ikke være mulig å utnytte det fulle potensialet hos de ansatte, fordi de mister motivasjon av mindre autonomi, som igjen vil skape negative konsekvenser for bedriftens lønnsomhet \parencite[s.~427]{einarsen2023}.

Innstramningen i hjemmekontorordningen vil kunne ha negative konsekvenser for de ansattes holdninger ovenfor ledelsen. Behren forklarer at de ansatte mener Hjemmekontorordningen fungerte godt, og at de fleste stiller seg uforstående til hvorfor ordningen må endres. Hun/han sier også at mange satt pris på fleksibiliteten ordningen gav, spesielt de ansatte med barn i barnehage \parencite{biernat2026}. Chris Argyris har kritisert teori X og mener at en slik organisering vil føre til at arbeidstakerne mister kontroll over arbeidsmiljø og blir en passiv del av organisasjonen \parencite[s.~427]{einarsen2023}. Det viser hvordan DNB ved en endring av hjemmekontorordningen kan være med på å gi de ansatte en følelse av mindre kontroll, noe som syntes gjennom hvordan Behren uttaler seg.

Det at de ansatte stiller seg uforstående ovenfor endringene og opplever at de mister fleksibilitet i hverdagen, kan få den konsekvensen at de opplever at ledelsen ikke lenger stoler på at de har en indre motivasjon for å arbeide. I stedet kan det oppstå en mistillit hvor det virker som ledelsen ikke stoler på arbeidsviljen til de ansatte, men opererer med et teori X menneskesyn i stedet.

Det er også mulig å bruke teori X til å forklare hvorfor innstrammingene ikke fungerer godt er at dersom mennesker motiveres av fritid, kan innstramminger rundt fleksibilitet og fritid skape mindre indre motivasjon.

\subsection{Teori Y og forutsetninger for indre motivasjon}

Teori Y er en motsats til teori X og påstår at mennesket er store ressurser i seg selv \parencite[s.~427]{einarsen2023}. Dette gjør at mennesket vil ha motivasjon til å gjennomføre sine arbeidsoppgaver dersom forholdene er til rettelagt \parencite[s.~427]{einarsen2023}. Mennesket vil til og med se på arbeidet som belønnende i seg selv på lik linje med lek og hvile. Teori-y tar med dette utgangspunkt i at mennesket vil verdsette varierte meningsfylte arbeidsoppgaver og anerkjennelse fra andre \parencite[s.~427]{einarsen2023}.

Ved benyttelse av denne teorien vil hjemmekontor derfor ikke være noe problem, fordi arbeidstakeren vil ha motivasjon til å gjennomføre sine oppgaver uavhengig av om det finnes et ytre press som \enquote{tvinger} deg til å være til stede på kontoret. Fordi mennesket motiveres av oppgaven i seg selv uavhengig av at det står noen og passer på at de gjør det de skal gjøre. Teorien presiserer at det er en forutsetning at det blir forholdene blir tilrettelagt for at arbeidstakeren kan løse sine arbeidsoppgaver.

Det kan argumenteres for at det å begrense muligheten for hjemmekontor vil føre til at forholdene da vil ligge bedre til rette for at arbeidsoppgavene kan gjennomføres. Dette kan videre kobles til DNB sitt argument om at økt tilstedeværelse på kontoret vil gjøre at arbeidsmiljøet vil bedres. DNB argumenterer videre for at endringen vil gjøre at arbeidstakerne vil kunne øke samarbeidet og læringen i organisasjonen, noe som viser hvordan DNB mener at omstendighetene vil ligge bedre til rette for at arbeidtakeren vil kunne gjennomføre sine oppgaver.

Det må videre understrekes at muligheten for hjemmekontor ikke forsvinner helt, men at det fortsatt vil være mulighet i dialog med leder. Dette viser hvordan DNB anerkjenner at individuelle tilpasninger fortsatt vil være nødvendig. Ved å koble dette til Y-teori kan det vise hvordan hjemmekontor fortsatt kan være den beste omstendigheten for en del av arbeiderne, men at DNB mener at det i hovedsak vil være bedre dersom flest mulig måter på kontoret.

Med tanke på Y-teori vil det derfor være sentralt hva som ansees som den beste tilretteleggelsen.

\subsection{Syntese: Motstridende perspektiver på kontorfellesskapet}

X og Y teori gir oss ulike syn på hvilken funksjon avviklingen av hjemmekontor vil ha. Der teori X argumenterer for at mennesket ikke sitter med indre motivasjon, men er grunnleggende late vil tilstedeværelse på kontoret være et nødvendig tiltak, mens teori Y vil se anderledels på dette siden mennesket har en indre motivasjon. Denne motivasjonen vil gjøre at arbeidstaker vil gjennomføre sine oppgaver selv om de ikke må være på kontoret ettersom de motiveres av oppgavene og ikke at de er på kontoret. Når det kommer til DNB sine argumenter vil teori x se på de som irrelevante fordi menneskets motivasjon ikke kommer fra arbeidsmiljø, samarbeid eller læring, men bare fra fritid og lønn. Teori Y vil på den andre siden trekke frem hvordan mennesket har indre motivasjon og omgivelsene er viktige for arbeidet. Dette passer med DNB sine argumenter om at omgivelsene blir bedre med den nye ordningen.
